\section{Conclusion}
\label{sec:conclusion}

Autoregressive and block-diffusion decoders do not sit at two points on a
latency scale. They have different cost \emph{shapes}: a line with no floor,
and a staircase with a floor that is mostly prefill (bounded above by
53--68\%). Almost everything a deployment needs to know follows from that
difference and from where the two shapes cross.

We measured both laws on one pinned stack, showed the staircase has sharp
edges exactly at canvas boundaries and uniform steps out to $L{=}2048$, found
the second variable that makes the diffusion law two-dimensional, and turned
the pair into a decision procedure whose constants all trace to committed
artifacts. Then we asked the question that motivated the work: does the
per-event advantage survive a fair comparison? At short output it does not.
The quality-matched autoregressive model is \footprintRatio$\times$ smaller,
batches \arBatchMult$\times$ better, and pays \inputSlopeRatio$\times$ less
per prompt token. At forced long output under a batch-1 constraint it does
survive, by $4.0$--$5.4\times$, and it brings two structural properties the
autoregressive model has no equivalent for: a per-request quality knob, and a
complete record at every point on the quality curve.

Two things we would emphasise to anyone extending this. First, the shape is
caused by the \emph{block}, not by the diffusion --- the denoising loop is at
least 32\% of the base cost (and likely closer to 47\%), so any decoder
emitting fixed-size blocks would inherit the same staircase, and the
interesting design space is block size rather than denoising schedule. Second, the four defects that biased our
own harness against the autoregressive baseline all lived in the
infrastructure around the model rather than in the model itself. That is not a
coincidence; it is what a comparison harness does when the person writing it
would prefer one of the two answers. The five checks in
\artifact{CHECKLIST.md} each take under five minutes, and each of them caught
something that had already changed a headline number.
